\documentclass{exam}
\usepackage{../commonheader}

\discnumber{2}
\title{Environment Diagrams and Lambda Expressions}
\date{September 16 -- September 19, 2025}

\begin{document}
\maketitle

\begin{meta}
    \textbf{Example Timeline}
    \begin{itemize}
        \item Intros + Ice Breaker [5 min]
            \subitem{Introduce yourself. Start off with an icebreaker of your choice!}
        \item Mini-lecture: Environment Diagrams [7 min]
            \subitem{Feel free to use the first question as an example/part of the mini-lecture.}
        \item Q1: FooBar [7 min]
        \item Q2: ABCD [10 min]
            \subitem{Feel free to spend more time on this if needed and skip the next question.}
        \item Q3: Swap [If Time]
        \item Mini-lecture: Lambda Expressions and Higher Order Functions [7 min]
        \item Q4: Def to Lambda [5 min]
        \item Q5: Make Interval [7 min]
        \item Q6: Lambda Environment Diagram [If Time]
            
    \end{itemize}
\end{meta}

\section{Environment Diagrams}
\begin{questions}
    \subimport{../../topics/environments/easy/}{foobar.tex}
    \begin{questionmeta}
        Suggested Time: 7 min
        \begin{itemize}
            \item Go over this question as an example for the mini-lecture, if needed.
            \item Make sure to explan when new frames are opened and how parent frames work when calling x and y at the end.
        \end{itemize}
    \end{questionmeta}

    \subimport{../../topics/functions-and-expressions/easy/env-diagram/}{abcd.tex}
    \begin{questionmeta}
        Suggested Time: 10 min
        \begin{itemize}
            \item The environment diagram questions can be a little long, so take extra time on this and skip Questions 3 and 6 if needed.
        \end{itemize}
    \end{questionmeta}
\end{questions}

\subimport{../../topics/functions-and-expressions/medium/env-diagram/}{swap.tex}
    \begin{questionmeta}
        Suggested Time: 10 min
        \begin{itemize}
            \item Optional, for if you have time after the first two questions or at the end.
        \end{itemize}
    \end{questionmeta}
\end{questions}


\section{Lambda Expressions and Higher Order Functions}
\begin{questions}
    \subimport{../../topics/hof/easy/}{def-to-lambda.tex}
    \begin{questionmeta}
        Suggested Time: 5 min
        \begin{itemize}
            \item This is a quick introduction to lambda functions. 
            \item Make sure to highlight the structure of lambda expressions and how they are called.
        \end{itemize}
    \end{questionmeta}

    \subimport{../../topics/hof/easy/}{make-interval.tex}
    \begin{questionmeta}
        Suggested Time:  7 min
    \end{questionmeta}

    \subimport{../../topics/environments/medium/}{abde.tex}
    \begin{questionmeta}
    Suggested Time: 7 min
    \begin{itemize}
            \item Optional, for if you have time. 
            \item This question is about how lambda expressions work in environment diagrams, so briefly go over how this works if you don't have time to complete the whole question.
        \end{itemize}
    \end{questionmeta}

\end{questions}

\end{document}